% BHU PYQ -- English: Drama-I (2025-26 NEP)
\documentclass[11pt,a4paper]{article}
\usepackage[a4paper,top=2.5cm,bottom=2.5cm,left=2.8cm,right=2.8cm,headheight=14pt]{geometry}
\usepackage{amsmath,amssymb}
\usepackage[T1]{fontenc}
\usepackage[utf8]{inputenc}
\usepackage{lmodern,microtype,fancyhdr,enumitem,hyperref,lastpage,parskip}

\hypersetup{colorlinks=true,urlcolor=black,linkcolor=black,pdftitle={ENGMJ32: Drama-I -- BHU NEP 2025-26}}

\pagestyle{fancy}\fancyhf{}
\renewcommand{\headrulewidth}{0.4pt}\renewcommand{\footrulewidth}{0.4pt}
\fancyhead[L]{\small \textit{Banaras Hindu University}}
\fancyhead[R]{\small \textit{ENGMJ32: Drama-I (2025-26)}}
\fancyfoot[C]{\small Page \thepage\ of \pageref{LastPage}}

\newlist{parts}{enumerate}{2}
\setlist[parts,1]{label=(\alph*),leftmargin=*,itemsep=4pt,topsep=3pt}
\setlist[parts,2]{label=(\roman*),leftmargin=*,itemsep=2pt,topsep=2pt}

\newcommand{\pts}[1]{\hfill \textit{[#1]}}

\begin{document}

\begin{center}
  {\large\bfseries BANARAS HINDU UNIVERSITY}\\[2pt]
  {\normalsize B. A. (Hons.) Semester III Examination, 2025-26 (NEP)}\\[6pt]
  \rule{0.8\linewidth}{0.5pt}\\[6pt]
  {\large\bfseries ENGLISH}\\[4pt]
  {\large\bfseries Paper Code: ENGMJ32 : Drama-I}\\[4pt]
  {\normalsize Time: 2:30 Hours \hfill Full Marks: 70}\\[2pt]
  {\small REF NO. 2025-26/D-III/084}
\end{center}
\rule{\linewidth}{0.4pt}
\smallskip

\noindent \textit{Instructions: Answer all Sections. Write your answers in your own words as far as practicable. Write your Roll No.\ at the top immediately on receipt of this question paper.}

\bigskip

\section*{SECTION -- A}
\noindent \textit{Note: Answer the following questions in about 400 words each:}\pts{3 $\times$ 12 = 36}

\begin{enumerate}[label=\textbf{\arabic*.},leftmargin=*,itemsep=12pt]

  \item Discuss the significance of equivocation in Shakespeare's play ``Macbeth''.\pts{12}

  \begin{center}\textbf{OR}\end{center}

  Discuss the role of Three Witches in ``Macbeth''.\pts{12}

  \item Discuss briefly the bargain made by Doctor Faustus. How does it portray the theme of the dangers of desire and temptation?\pts{12}

  \begin{center}\textbf{OR}\end{center}

  Give a character sketch of Mephistopheles in Christopher Marlowe's ``Doctor Faustus''.\pts{12}

  \item How does ``The Merchant of Venice'' explore the themes of prejudice, justice, and mercy?\pts{12}

  \begin{center}\textbf{OR}\end{center}

  Write a note on the female characters in ``The Merchant of Venice''. What ideas does it give you about Renaissance society and gender?\pts{12}

\end{enumerate}

\bigskip

\section*{SECTION -- B}
\noindent \textit{Note: Answer the following questions in about 150 words each:}\pts{4 $\times$ 6 = 24}

\begin{enumerate}[label=\textbf{\arabic*.},leftmargin=*,start=4,itemsep=10pt]

  \item Write a brief note on Morality Plays.\pts{6}

  \begin{center}\textbf{OR}\end{center}

  Write a brief note on some major elements of Drama.\pts{6}

  \item Answer with reference to context:\pts{6}
  \begin{parts}
    \item ``Out, damned spot! out, I say! One; two: why, then 'tis time to do't. Hell is murky!'' \textit{(Macbeth)}
    \item ``Was this the face that launched a thousand ships,\\And burnt the topless towers of Ilium?'' \textit{(Doctor Faustus)}
  \end{parts}

  \item Answer with reference to context:\pts{6}
  \begin{parts}
    \item ``Thou know'st that all my fortunes are at sea;\\Neither have I money, nor commodity\\To raise a present sum: therefore go forth;\\Try what my credit can in Venice do:'' \textit{(The Merchant of Venice)}
    \item ``Hath a dog money? Is it possible\\A cur can lend three thousand ducats?'' \textit{(The Merchant of Venice)}
  \end{parts}

  \item Write a brief note on Doctor Faustus as a Renaissance man.\pts{6}

  \begin{center}\textbf{OR}\end{center}

  Write a short note on Soliloquy in Shakespearean tragedy.\pts{6}

\end{enumerate}

\bigskip

\section*{SECTION -- C}
\noindent \textit{Note: Short Answer Type Questions. Answer any \textbf{five} of the following questions in about 50 words each:}\pts{5 $\times$ 2 = 10}

\begin{enumerate}[label=(\alph*),leftmargin=*,itemsep=6pt]

  \item State the reasons why Shylock and Antonio despised each other.\pts{2}
  \item Discuss two classical allusions from ``The Merchant of Venice''.\pts{2}
  \item How does Mephistopheles fulfill Faustus's demand to meet Helen of Troy?\pts{2}
  \item State the difference between a Morality Play and a Mystery Play.\pts{2}
  \item Explain two imageries used in ``Macbeth''.\pts{2}
  \item What is meant by exposition of the play? Explain with an example.\pts{2}

\end{enumerate}

\end{document}
